\section{Fixed one-upward Pachner descent before normal search}
\label{sec:pachner-descent-crushing}

A triangulation with no accepted strict $3\to2$ move can still admit a
short simplification that first inserts one tetrahedron.  The optional
policy in this section searches completely for a first strict descent
with at most one total $2\to3$ move, accepts an independently replayed
geometric macro, and restarts normal search on its smaller endpoint.
The intermediate size can increase, but each accepted macro decreases
the source size by exactly one.  This gives polynomial preprocessing
without a monotonicity assumption about the normal-surface cones.

The source entry point is \path{reduce_pachner_descent_once}; the direct
and factored crushing interfaces expose it through
\path{reduce_pachner_descent=True}, with default false.  The underlying
bounded-descent and cover primitives are those of
Sections~\ref{sec:pachner-native} and~\ref{sec:cover-oracle}.
Unlike the transported-disc epoch policy of
Section~\ref{sec:pachner-epochs}, this adapter uses only zero heights
and does not query for a normal disc.  Its accepted evidence certifies
a change of triangulation of the same compact manifold.

\subsection{A first descent has one or three elementary moves}

Write $U$ for a $2\to3$ expansion and $D$ for a $3\to2$ contraction.
Here the allowance is one \emph{total} expansion in one call, rather
than a bound on the maximum excess of tetrahedra.  Starting with $t$
tetrahedra, a prefix with $u$ expansions and $d$ contractions has
$t+u-d$ tetrahedra.  Before a first strict decrease,
\[
                         d\le u\le1.
\]
Consequently the only possible nonterminal words are the empty word,
$U$, and $UD$.  A completed first strict descent is either $D$ or $UDD$.
Its endpoint has $t-1$ tetrahedra and all its intermediate sources have
at most $t+1$ tetrahedra.  A packet with an earlier strict decrease
followed by further moves is outside this fixed macro policy, even if
its eventual endpoint is smaller and every individual move is valid.

The search calls \path{find_pachner_descent} with upward allowance one,
the sleep-set method, the exact oracle cover backend, and all node,
region, cooperative-work and oracle-state caps removed.  The default
connected original-cover radius is
\[
                              R=\min(t,6).
\]
The first-descent locality theorem recalled in
Section~\ref{sec:pachner-native} supplies this sufficient radius:
if a formal descent using at most one expansion exists, a connected
first-descent witness consumes at most six original tetrahedra.
This statement concerns existence of an appropriate witness; it does
not assert that the original consumption of every trace is connected.
The cover oracle admits a footprint precisely when some such connected
cover contains it.  Its current witness does not restrict future
eligibility; later queries use the consumed footprint itself.

The endpoint predicate stops at the first visited smaller source.
Complete exhaustion returns \path{UNCHANGED}; success returns
\path{REDUCED} and a proof packet.  There is no finite internal search
allowance that can be mistaken for exhaustion in this adapter.
A caller interruption propagates instead of returning either outcome.
An unchanged result concerns this bounded formal move family only.
It proves neither a minimum triangulation nor a knot verdict; in
particular it is not a separately replayable negative normal proof.

\subsection{Relative-boundary topology and the zero cochain}

Both elementary moves exchange the two formal triangulations of the
triangular bipyramid used in
Proposition~\ref{prop:strict-pachner32-relative}.  Their common six
formal boundary triangles carry the same complete attachment maps.
The homeomorphism fixing this formal boundary descends through all
those relations, including attachments between two boundary facets of
the same region.  The inverse replacement satisfies the same argument.
Thus a composition of independently verified elementary moves preserves
the underlying compact manifold and its boundary, even when the formal
region has identified boundary vertices, edges or facets in the quotient.

Each original, intermediate and final source is separately checked as
a finite connected orientable compact triangulated manifold with exactly
one torus boundary.  The accepted region tetrahedra are distinct and
the internal facet maps, survivor records and formal boundary maps are
checked, not inferred from matching source sizes or homology.  The
source API permits any number of global vertices.  Since neither
formal triangulation introduces an interior vertex of the bipyramid,
these replacements preserve the ambient vertex equivalence classes.
In particular a one-vertex source remains one-vertex.

For the existing transport machinery, give each tetrahedron the height
row $(0,0,0,0)$.  All signed edge differences vanish, so these rows
define a consistent integral cochain.  On a formal move region, its
connected edge graph forces all five aligned bipyramid heights to be
equal.  Normalization at the first formal vertex makes all five zero.
The replacement rows, being restrictions of these values, are again
zero; surviving rows stay zero as well.  Induction proves this for
every intermediate source.  The associated seven local normal
coordinates vanish in every tetrahedron, as do the Euler and normal-disc
jumps.  No coordinate size growth occurs.

This zero cochain carries the empty normal surface.  It supplies a
uniform way to reuse exact geometric transport and replay, but provides
no positive Euler ray, essential disc, peripheral generator or normal
surface identification between the old and new cones.  The adapter
does not choose a cocycle seed, run an LP, or invoke disc discovery.
The proof explicitly requires the endpoint's disc payload to be absent.

\subsection{Source binding and independent strict acceptance}

The new outer packet has exactly four fields:
\path{schema}, \path{source_sha256}, \path{certificate}, and
\path{successor}.  The schema value is
\path{normal-pachner-descent-v1}.  The digest is SHA-256 of the complete
labelled source encoded as canonical JSON with sorted keys and compact
separators.  The successor is its full labelled face table, and the
nested certificate uses the existing
\path{pachner-cochain-endpoint-v1} schema.

The independent entry point \path{verify_normal_pachner_descent}
receives the original source explicitly.  It reconstructs its zero
height rows, checks the digest and exact field sets, and independently
replays every geometry and transport record through the existing
endpoint and cover inspectors.  This establishes each relative-boundary
replacement, every intermediate finite source, and a connected original
cover of size at most $\min(t,6)$ containing all consumed original cells.
The digest identifies the labelled source; actual replay from that
source establishes the mathematical claim.

The new verifier also enforces conditions beyond the legacy endpoint
inspector's eventual strict decrease:
\begin{enumerate}
\item The advertised upward allowance is the literal integer one; the
      trace has between one and three moves and at most one actual
      expansion.
\item Every proper prefix has at least $t$ tetrahedra, and the nonempty
      final source has exactly $t-1$ tetrahedra.  Independent move replay
      checks the one-tetrahedron change at each step.
\item Every recorded height row, seven-coordinate row, five-entry
      formal bipyramid height row and both jumps consists of literal
      integer zeros; the final coordinate rows are zero and the
      \path{disc_certificate} value is \path{None}.
\item The declared successor agrees with the independently replayed
      endpoint by type-sensitive recursive equality.  Boolean or
      floating-point replacements for integer labels are not accepted
      as equal source data.
\end{enumerate}
The resulting accepted traces are exactly of the shapes $D$ or $UDD$.
The independent checker imports replay code, not the descent producer;
it needs no successful rerun of discovery.

Both source adapter and checker guard the caller's cancellation callback
across nested legacy interfaces.  An internal carrier preserves the
original exception object and prevents a caller-raised malformed-input
exception or search-limit exception from becoming an ordinary rejection
or exhaustion.  Malformed evidence is rejected; an interrupted accepted
packet is not silently appended to the outer chain.  Complexity bounds
below count the implementation's bookkeeping and charge arbitrary work
performed by a caller-supplied callback separately.

\subsection{Polynomial cost of the actual fixed search}

The finite upward allowance already gives a direct bound for this
particular implementation.  The argument includes exact marked keys,
cover caches and lazy transition construction, rather than treating
either the cover oracle or associative lookup as a free operation.
Let $L=\log(t+2)$ and let $P(s)$ bound the local work at size $s$:
complete source preparation, legal-event enumeration, one elementary
construction and zero transport, exact serialization and independent
local replay.  As in Section~\ref{sec:strict-pachner32}, the conservative
choice
\[
                           P(s)=O(s^2\log^2(s+2))
\]
covers these operations on explicit dense sources.  It permits linear
scans for local associative lookups.  It does \emph{not} yet include
lookups in the global caches accumulated over the whole search.

\begin{proposition}[Fixed-policy search and proof bounds]
\label{prop:pachner-descent-crushing-cost}
For a validated $t$-tetrahedron source, one completed call to the fixed
zero-height adapter has the conservative bit bounds
\[
 \text{discovery time }O(t^6L^2),\qquad
 \text{working space }O(t^3L).
\]
An accepted macro packet has $O(tL)$ bits and can be independently
replayed in $O(t^2L^2)$ bit operations.  These bounds require no
average-case hash-table assumption.
\end{proposition}
\begin{proof}
At a live source with $s\le t+1$ tetrahedra there are $O(t)$ eligible
events.  Upward events arise from paired facets between distinct
tetrahedra; there are at most $2s$ such pairs.  Downward events arise
from degree-three global edge classes using three distinct tetrahedra;
there are $O(s)$ classes among the $6s$ local edge incidences.
Preparing the full table and checking the formal regions is covered by
$P(s)$; this event count does not declare every combinatorial edge
eligible without checking its actual region.

Before success, only the empty, $U$ and $UD$ prefixes can be expanded.
Thus there are
\[
                  N=O(t^2)\quad\hbox{visited frames},\qquad
                  Q=O(t^3)\quad\hbox{candidate attempts},
\]
with at most one additional selected terminal frame.  The code
materializes a queued move through \path{_advance} only when that frame
is popped.  Each visited nonterminal frame generates its event list
through \path{_events}; it need not materialize all candidate children.
Sleep pruning and failed cover queries can only remove these paths.
Duplicate marked endpoints do not prune their continuation frames, so
the bound includes such frames as well.

The smaller-source predicate is queried before generating children at
the first occurrence of each exact marked endpoint.  Its key contains
one record per live tetrahedron.  Consequently a smaller source cannot
be concealed by a key of a nonsmaller source; the first visit to any
smaller key stops the entire query.  Deduplication therefore cannot
create longer uncounted paths below the first-descent threshold.

The cover graph has maximum degree four and $R\le6$.  Each footprint
query asks for a rooted completion inside a set of size at most six.
At any stage its frontier has size at most $24$, and extension depth
is at most six.  Exhaustive canonical-parent completion therefore has
an absolute constant number of local extension possibilities, with
$O(L)$-bit original labels.  Oversized footprints are rejected directly.
A queried footprint contains at most eight originals along these
three-event traces, so its cache key also has constant cardinality.
The selected cover witness is a constant-size record, not an ambient
index over all covers.

Across the query, the number of interned incarnation identifiers is
$O(t+N)$, since each materialized move creates at most three.  Their
bit lengths are $O(L)$.  Interning keys contain an event's fixed-arity
integer incarnation ports and a birth slot; they do not recursively
expand the entire ancestry history.  A marked endpoint key has $O(t)$
face and height records and thus $O(tL)$ bits.  There are at most $N$
such keys and at most $Q$ completed footprint-cache records.  These
require $O(NtL+QL)=O(t^3L)$ bits.  The depth-three DFS stack, event lists,
sleep payloads and constant-length full-source traces fit the same
bound; even retaining a full source with every visited frame would fit
it.  Unvisited siblings are stored as parent/event records.  Temporary
certificate copies and final replay do not increase this order.

For a deterministic worst-case time bound, charge a full scan of the
global associative table for every lookup.  The $Q$ cover-cache
lookups cost at most $O(Q^2L)$, since their keys have constant
cardinality.  Marked-key membership costs at most $O(N^2tL)$.
Fixed-arity incarnation interning is no larger.  Comparing sleep
obligations across all candidate events costs at most $O(Nt^2L)$.
All materialized geometry and local preparation cost $O(NP(t+1))$.
In particular the sum
\[
 O\bigl(NP(t+1)+Q^2L+N^2tL+Nt^2L\bigr)
                \subseteq O(t^6L^2)
\]
covers global collision scans separately from the local bound $P$.
Source hashing, proof construction and the adapter's independent final
replay fit this conservative estimate.

Finally, an accepted proof has at most three full labelled successor
snapshots.  Each has $O(tL)$ bits; its zero height and coordinate rows
have $O(t)$ entries, and its formal move and cover records have
$O(L)$ bits.  Repeating the final source in the outer successor field
changes only a constant factor.  Constantly many source checks,
transport replays and literal comparisons cost $O(P(t+1))$, proving
the packet and replay bounds.
\end{proof}

The callback counter \path{search_work} is the legacy descent call's
\path{aggregate_work}; it counts cooperative checkpoints through that
call's preparation, search and legacy positive replay.  It is not a
bit-operation counter or a timing measurement.  It does not include
all outer-adapter hashing, source checks or subsequent driver replay.
Likewise \path{search_nodes} counts popped search frames, rather than
candidate events, normal-search nodes or cover completions.  The
asymptotic proof charges those additional operations explicitly.

Restarting this preprocessing alone accepts fewer than $t$ macros.
Its full discovery cost is at most $O(t^7L^2)$, peak working space is
$O(t^3L)$ when completed query caches are released, and the retained
accepted packet chain uses $O(t^2L)$ bits.  It contains at most
$3(t-1)$ elementary moves and has intermediate size at most $t+1$.
Replaying this whole accepted chain costs at most $O(t^3L^2)$.
These are conservative bounds for fixed geometric preprocessing,
separate from all subsequent Euler decisions and normal-ray extraction.

\subsection{Mixed-chain scheduling, budgets and source resets}

At each current source the crushing driver first attempts enabled
vertex reduction, then inverse boundary layering, then the strict
$3\to2$ scan, and finally this one-upward macro.  A successful
preprocessor restarts the loop immediately.  Only after these enabled
operations return unchanged does normal Euler search prepare its model.
After a macro, incidence tables, source digests, matching equations,
LP models, dual certificates and learned conflicts belong to the new
source and are reconstructed.  No old normal-search cache is transported
by the zero-height geometric certificate.

The mixed chain uses \path{diagram-crushing-v4} whenever this option
is enabled.  It retains the same exact outer fields as the earlier
chain versions and accepts normal crushes, inverse boundary layers,
strict $3\to2$ packets and the new macro packets.  Versions one, two
and three retain their earlier accepted step kinds and reject the new
macro.  The independent chain checker authenticates the diagram
exterior and any initial shellings before replaying every step from
its actual predecessor.  It requires strict size decrease between
outer steps.  The intermediate expansion is inside one verified macro,
whose endpoint satisfies this same decreasing measure.

The shared \path{max_crushes} allowance counts accepted outer steps:
\[
 \text{normal crushes}+\text{inverse layers}
 +\text{strict }3\to2\text{ moves}+\text{one-upward macros}.
\]
A $UDD$ macro uses one unit, while its one upward and two downward
elementary moves are recorded separately.  Under the common
\path{pachner_descent_} field prefix, \path{calls}, \path{nodes} and
\path{work} accumulate completed discovery calls.  The suffixes
\path{macros}, \path{upward_moves} and \path{downward_moves} instead
record accepted progress.
Discovery counters are accumulated before testing whether a discovered
macro may be accepted under the shared allowance.  Thus a successful
search denied acceptance still contributes its reported discovery work,
while the current source and accepted-move counters remain unchanged.
The reduction budget does not bound this polynomial discovery work.
An exception during an outstanding call propagates and yields no
completed-call statistics to claim as complete work.

The source producer independently verifies a successful packet, and
the driver replays it again before appending it.  If the allowance is
exhausted, the driver returns inconclusive with its earlier authenticated
partial chain.  A nonempty geometric successor alone proves neither
triviality nor nontriviality.  Recognition still requires the usual
terminal normal crushing proof or complete negative normal certificate.
The existing outer recognition schemas use the new chain through the
same independent chain verifier.

The factored interface passes this policy to each attempted visible
factor and charges accepted macros once against the same global
reduction allowance shared with the other step kinds.  Additive search
statistics include every attempted child.  Its proof quorum remains
one verified knotted child or verified unknot proofs for every child;
a geometric macro does not itself satisfy either quorum.

\subsection{Native validation scope}

The retained regression run passes $336$ tests across $38$ modules with
Regina available and no skipped tests.  It covers the exact LP and normal
crushing paths, legacy cover and cochain interfaces, and the new producer,
independent checker and mixed-chain integration.  Controls include both
single contractions and three-move first descents, several global vertices,
formal-region selfattachments, caller cancellation through nested replay,
and shared budgets across incomplete factor attempts.  A genuine complete
trefoil call and its factorwise composition verify through the new chain.

A final focused seven-test verifier rerun strengthens the outer successor
control: its corrupted boolean or floating-point label is rejected while
the unchanged inner geometric proof still independently passes.  A
geometrically valid $DUD$ trace also passes the legacy eventual-descent
inspector but is rejected by the new first-descent policy.  These controls
check the stricter interface contract, rather than treating geometric
reachability alone as proof of the fixed-policy bound.  Native code stayed
unchanged between the broad and focused runs.  The commands and logs are
retained under \path{data/pachner-descent-*}; the core publication helper
pins the implementation, fixture dependencies and rebuilt article.
These tests establish finite regression evidence, not a scaling theorem
or a measurement of recognition speed.

\subsection{What this gives toward the complexity objective}

Every accepted outer step strictly decreases tetrahedron count, so
interleaving this policy with the other strict reductions still permits
only $O(t)$ source stages and hence $O(t)$ fixed-macro queries.  The
polynomial preprocessing bound is therefore compatible with the
factorwise cost accounting of Section~\ref{sec:factored-crushing}.
It also preserves the restricted visible-factor-size corollary there:
the explicit complete tree policy has exponential work in the largest
child's linear-size source, and this added polynomial term does not
change that bound.

There is no proved inequality $h(T')\le h(T)$ for a returned macro,
nor a proved transport of the positive formal root rays or their
conflict sets.  Even the primitive negative-source resolving parameter
is defined using the current matching cone and canonical Euler data,
which are changed by retriangulation.  The relative-boundary
homeomorphism and the transported empty surface do not identify these
objects.  An increase is not ruled out by this proof, and no concrete
increase witness is asserted here.  Fewer tetrahedra alone therefore
does not certify fewer LP nodes or a faster complete recognition run.

Conditional bounds involving a polylogarithmic resolving parameter
must still control that parameter on \emph{every} authenticated source
at which normal search actually runs, including the new macro endpoints.
The fixed search supplies no such uniform control.  A source stationary
under this total-one-upward family need not be globally minimal or
cheap to recognize.  Thus the policy establishes a certified polynomial
preprocessing operation; it supplies no general quasi-polynomial bound
for unknot recognition.

\subsection{Archived staged experiments and exact endpoint parameters}
\label{sec:pachner-descent-crushing-staged}

The package \path{fast/pachner_descent_research} preserves two distinct
historical evidence groups.  Its manifest records hashes and lengths of
both stored files and their uncompressed contents; embedded native hashes
identify the code used when an old artifact was generated.  The current
independent checker replays those artifacts with the current pinned
verifiers.  This is fresh proof replay, without fresh LP discovery or a
new Regina oracle run.  Historical generation pins are not represented
as current native source pins.

The earlier group, \path{strict-descent-next-design}, includes the exact
strict-$3\to2$-preprocessed sources for the braid closures
\[
 \begin{split}
 T_{25}&=\widehat{\sigma_1^5}\quad\text{on two strands},\\
 T_{27}&=\widehat{\sigma_1^7}\quad\text{on two strands},\\
 W&=\widehat{(\sigma_1\sigma_2^{-1})^4}
                         \quad\text{on three strands}.
 \end{split}
\]
Their respective source sizes are $15,21,34$.  The explicit braid data
and authenticated native exterior chains define the experiment.  On the
15-tetrahedron $T_{25}$ source, a complete negative tree with split set
$\{0,1,4\}$ proves $h\le3$.  Three independently checked positive
formal root rays have pairwise disjoint conflicts
\[
             \{0,7\},\qquad\{1,11,12\},\qquad\{4,10,13\},
\]
and prove the matching lower bound.  Thus $h=3$ on that labelled source.
On the 21-tetrahedron $T_{27}$ source, three verified rays have disjoint
conflicts $\{0,7\}$, $\{1,10,16,17\}$ and $\{4,6\}$, proving $h\ge3$.
Its capped full-root query is inconclusive, so there is no corresponding
upper bound from that query.  The 34-tetrahedron weaving source also
remains inconclusive.  Across this group, eleven distinct exposed rays
are checked with the original STANDARD matching equations and exact
rank; no new ray enumeration is needed for replay.

That group also retains three separately capped supplied-source
one-upward searches from the older $8$-tetrahedron trefoil,
$12$-tetrahedron figure-eight and $17$-tetrahedron two-trefoil sources.
They produce first-descent endpoints of sizes $7,11,16$, each by $UDD$.
These local packets predate the new source adapter and do not by
themselves supply complete recognition outcomes or endpoint bounds on $h$.

The second group, \path{fixed-u1-staged-audit}, records exactly twenty
completed calls: two whole-driver circle calls; twelve preprocessing-only
whole-driver calls on six noncircle diagrams, with the new policy off
and on; and six separate Euler queries on the enabled endpoints.
The twelve preprocessing calls set the normal tree-node allowance to
zero and perform zero LP calls.  Every baseline chain is literally a
prefix of its enabled chain, with identical original exterior and
shelling data.  The six later source queries use the conflict-based
transversal policy, polynomial LP backend, $256$ nodes and $10{,}000$
each of pivots, anchors, candidate subsets, ellipsoid iterations and
feasibility calls.

In the following table, $t_0$ and $t_1$ are the baseline and enabled
preprocessing endpoint sizes; $M$ counts accepted macros.  The node and
pivot columns belong only to the separate supplied-source Euler query.
Here $T$ is the trefoil, $F$ the figure eight, and $TT$ is the closure
of the three-strand braid $\sigma_1^3\sigma_2^3$.  ``Negative'' means
a complete full-root negative normal proof, while ``capped'' records
an inconclusive return.
\begin{center}
\begin{tabular}{lrrrrrlc}
\hline
Input & $t_0$ & $t_1$ & $M$ & Nodes & Pivots & Query & $h(T_1)$\\
\hline
$T$      & 8  & 7  & 1 & 5  & 259   & negative & $1$\\
$F$      & 12 & 10 & 1 & 1  & 676   & negative & $0$\\
$TT$     & 17 & 16 & 1 & 16 & 3526  & negative & $2$\\
$T_{25}$ & 15 & 11 & 3 & 7  & 693   & negative & $1$\\
$T_{27}$ & 21 & 17 & 4 & 15 & 2545  & negative & $2$\\
$W$      & 34 & 31 & 3 & 1  & 10000 & capped   & unknown\\
\hline
\end{tabular}
\end{center}
All thirteen accepted macros in this staged group have one upward and
two downward moves.  Strict $3\to2$ preprocessing runs again after
each macro: in particular the figure-eight sequence goes from $12$ to
$11$ by the macro and then to $10$ by a strict contraction.  The
$T_{25}$ chain also has an additional strict contraction beyond its
three accepted macros.  The other enabled preprocessing rules remain
part of these actual chains.

Each of the five negative source proofs is combined with its exact
authenticated preprocessing chain to form an independently verified
knottedness certificate.  These are assembled proofs, not five new
whole-driver discovery calls.  In particular summing their separate
source-query counters with preprocessing counters would not create a
measurement of a fresh recognition run.  Both circle whole calls prove
the unknot.  The weaving case retains its partial chain and capped
source query; it supplies neither a complete negative proof nor a
nontrivial bound on $h$.

\paragraph{Why the endpoint values of $h$ are exact.}
The five complete negative trees split only at the following sets,
in the table's order:
\[
        \{2\},\qquad\varnothing,\qquad\{0,10\},\qquad
        \{7\},\qquad\{10,14\}.
\]
Proposition~\ref{prop:strict-pachner32-tree-upper-bound} therefore
gives upper bounds $1,0,2,1,2$.  Its argument applies to the entire
positive formal root-ray family: a ray with no conflict in the split
set could follow a compatible branch to a leaf, contradicting that
leaf's independently checked exclusion of positive Euler vectors.
This upper-bound argument does not use completeness of a finite
sample of exposed rays.

Matching lower bounds use already exposed rays.  The following selected
conflicts suffice; $(c,r)$ denotes the number of retained STANDARD
columns and their exact rank.
\begin{center}
\begin{tabular}{llc}
\hline
Labelled endpoint & Selected conflicts & $(c,r)$\\
\hline
$T$, $7$ tetrahedra & $\{2,6\}$ & $(35,33)$\\
$TT$, $16$ tetrahedra & $\{0,4\}$, $\{10,14\}$
                                      & $(73,71)$, $(77,75)$\\
$T_{25}$, $11$ tetrahedra & $\{2,7\}$ & $(55,53)$\\
$T_{27}$, $17$ tetrahedra & $\{2,14\}$, $\{0,10,11\}$
                                      & $(85,83)$, $(85,83)$\\
\hline
\end{tabular}
\end{center}
Every listed ray has canonical Euler value one.  The independent
calculation checks nonnegative integral coordinates, original matching,
primitive quadrilateral scaling, canonical triangle normalization and
incidence Euler.  Each source has one finite vertex and two boundary
triangles.  All $4t$ triangle columns are retained, with rank
$4t-1$, together with the positive quadrilateral support.  Thus the
triangle-only kernel has dimension one and the projected quadrilateral
support nullity is
\[
                             c-r-1=1.
\]
This verifies extremality in the unrestricted quadrilateral root cone;
it does not mistake extremality in a discovery mask for root extremality.
A nonempty conflict needs at least one hitting tetrahedron, and each
displayed disjoint pair needs at least two.  These lower bounds equal
the complete-tree upper bounds.  The checker replays all ten distinct
exposed rays in the staged archive, including those not needed in the
displayed witness selection.

For the figure-eight endpoint, $h=0$ follows directly from the negative
unrestricted-root leaf.  It does not follow from the absence of exposed
rays in a sample.  The older 12-tetrahedron source had the proved upper
bound $h\le1$; that earlier evidence did not establish $h=1$, so this
comparison alone is not a proved strict decrease of $h$.

There are nevertheless two proved strict source-specific decreases:
$T_{25}$ changes from exact $h=3$ at $15$ tetrahedra to exact $h=1$ at
$11$, and $T_{27}$ changes from $h\ge3$ at $21$ tetrahedra to exact
$h=2$ at $17$.  All source identifications use full labelled arrays
and hashes, not knot type or tetrahedron count alone.  These finite
endpoint comparisons establish neither monotonicity at each intervening
move nor a uniform bound after exhaustion of the fixed one-upward
family.  Discovery counters in this archive are retained nonnegative
counts checked against configured caps; independent replay verifies
the mathematical packets and accepted-move totals, not an execution
transcript of every LP operation.

\subsection{Fourteen fresh complete-driver audit calls}
\label{sec:pachner-descent-crushing-audit}

The separate \path{data/pachner-descent-audit.json} contains fourteen
fresh calls of the raw whole-diagram crushing recognizer: the circle
and the six noncircle inputs above, each with the macro disabled and
enabled.  The visible-factor wrapper is not invoked, including on $TT$.
These calls actually run preprocessing, normal search and final
recognition replay together.  They are distinct from the twenty-call
staged protocol and its five assembled proofs.

Both arms enable vertex reduction, boundary peeling and the strict
$3\to2$ scan.  They use the source-model transversal search with
conflict-based candidates, the envelope strategy and the polynomial
LP policy with its default bounded fast-path prefix
(\path{lp_fast_path_pivots=None}).  The shared limits
are $256$ normal tree nodes and $10{,}000$ each of pivots, anchors,
candidate hitting subsets, ellipsoid iterations and feasibility calls.
The option map also fixes $256$ sectors, $10{,}000$ bases, diversified
sector order and quantum $64$.  The reduction allowance is unlimited;
strict source descent still bounds accepted outer operations.  Only
the macro-enable flag differs between the arms.  There is no wall
cutoff or removal of an unsuccessful case.

All fourteen calls returned, with no exception or unreturned call.
The audit independently replays every conclusive proof or partial
chain and retains every complete returned result before replay.
The native and research dependency map has $406$ unchanged file pins
through the run.  The table records reported normal-search nodes $N$
and pivots $P$, independently recoverable normal crushes $C$, strict
$3\to2$ moves $S$, accepted macros $M$, and final tetrahedron count
$t_f$.  Here U, K and I denote \path{UNKNOT}, \path{KNOTTED} and
\path{INCONCLUSIVE}; a dash in the last column means the terminal
empty crushing outcome, not an unrecorded source.
\begin{center}\small
\begin{tabular}{llcrrrrrr}
\toprule
Input & Macro & Result & $N$ & $P$ & $C$ & $S$ & $M$ & $t_f$\\
\midrule
Circle & off & U & 0 & 0 & 1 & 0 & 0 & ---\\
Circle & on  & U & 0 & 0 & 1 & 0 & 0 & ---\\
$T$ & off & K & 7 & 428 & 12 & 2 & 0 & 8\\
$T$ & on  & K & 5 & 259 & 12 & 2 & 1 & 7\\
$F$ & off & K & 7 & 1917 & 16 & 1 & 0 & 12\\
$F$ & on  & K & 1 & 676 & 16 & 2 & 1 & 10\\
$TT$ & off & K & 56 & 3398 & 24 & 6 & 0 & 17\\
$TT$ & on  & K & 16 & 3526 & 24 & 6 & 1 & 16\\
$T_{25}$ & off & K & 27 & 7113 & 20 & 5 & 0 & 15\\
$T_{25}$ & on  & K & 7 & 693 & 20 & 6 & 3 & 11\\
$T_{27}$ & off & I & 21 & 10000 & 28 & 11 & 0 & 21\\
$T_{27}$ & on  & K & 15 & 2545 & 28 & 11 & 4 & 17\\
$W$ & off & I & 1 & 10000 & 32 & 3 & 0 & 34\\
$W$ & on  & I & 1 & 10000 & 32 & 3 & 3 & 31\\
\bottomrule
\end{tabular}
\end{center}
All boundary-unlayering counts and both LP fallback work counts are
zero in these calls.  The enabled noncircle runs make the following
macro-discovery calls, including their final unsuccessful call:
\begin{center}\small
\begin{tabular}{lrrr}
\toprule
Input & Calls & Popped macro frames & Cooperative checkpoints\\
\midrule
$T$ & 2 & 41 & 33259\\
$F$ & 2 & 68 & 76728\\
$TT$ & 2 & 129 & 215368\\
$T_{25}$ & 4 & 125 & 160180\\
$T_{27}$ & 5 & 168 & 310361\\
$W$ & 4 & 202 & 628206\\
\bottomrule
\end{tabular}
\end{center}
The circle terminates through vertex-reduction crushing before any
macro search.  Macro frames and checkpoints are separate from $N$ and
$P$ and have the diagnostic scope described in the cost proof above.
They retain the preprocessing effort instead of treating a smaller
normal query as a free transformation.

For $T_{27}$ the enabled policy completes a knottedness proof within
the declared $10{,}000$-pivot limit where the baseline returns only
a partial chain.  This is a finite improvement in completion under
the same allowances, not an elapsed-time ratio between equal complete
tasks.  Both weaving calls reach that pivot limit and remain
inconclusive despite the source reduction from $34$ to $31$
tetrahedra.  On $TT$, nodes decrease from $56$ to $16$ while pivots
increase from $3398$ to $3526$.  Thus even these retained controls
do not support a uniform decrease in all normal-search work measures.
The public macro default remains disabled.

\paragraph{Binding the fresh sources to the archived parameter proofs.}
The frozen package checker verifies the whole-call proofs and archive
proofs separately.  The additional helper
\path{data/pachner_descent_endpoint_binding.py} checks their relationship
explicitly.  Its retained output,
\path{data/pachner-descent-endpoint-bindings.json}, verifies all twelve
noncircle endpoint comparisons by exact type-sensitive equality of
their full labelled source tables.  Each also passes the full finite
source validator and has one vertex and two boundary triangles.
The helper binds the input PDs, option maps, source pins and complete
coverage; its source-equality checks supplement, rather than replace,
the package's independent replay of every complete or partial proof.
Consequently the staged endpoint values and bounds on $h$ apply to
these same fresh sources.  This uses equality of sources, not a
theorem transporting $h$ through an arbitrary homeomorphism.

Separately, the retained corruption controls reject $49$ mutations
of saved evidence, with their unmodified controls accepted.  Their
scope includes genuine circle and macro packets, option and counter
contracts, audit coverage, and source binding.  One control supplies
a different labelled trefoil endpoint with a valid native recognition
proof: ordinary native result replay accepts that proof, while the
explicit archive-endpoint binding rejects the false equality claim.
This distinguishes a correct knot verdict from the additional source
identity needed to publish a particular value of $h$.  The control
script, fixture and output are retained under
\path{data/pachner-descent-corruption-*}.

Accepted move totals and final source sizes follow from replayed
packets.  Search statistics satisfy literal nonnegative-integer,
configured-cap and aggregation checks, but are reported counters rather
than an independently reconstructed LP execution history.  These
fourteen fresh calls and exact endpoint checks establish finite
behavior and proof provenance; they provide no asymptotic recognition
bound or general monotonicity theorem.

\subsection{Fifty paired complete-call timing controls}
\label{sec:pachner-descent-crushing-timing}

The predeclared benchmark in
\path{data/pachner-descent-benchmark.json} completes all fifty calls
on circle, trefoil, figure eight, $TT$ and $T_{25}$.  Each input has
one warmup per arm, three pairs with alternating arm order, and one
baseline A/A pair.  There is no wall cutoff.  The $T_{27}$ and weaving
inputs remain in the fourteen-call audit above but are outside this
predeclared timing set; their capped baseline outcomes are not used
to form recognition-speed ratios.

Each timer includes fresh whole-driver recognition, independent replay
of its conclusive proof or partial chain, canonical serialization of
the entire returned result and its SHA-256 hash.  Module initialization,
prepared input PDs and reference-result serialization are outside the
timer.  A runtime check confirms that no additional native module is
initialized inside a measured call.  Warmup results are retained in
full; every sample's result hash and serialized length agree with its
complete audited arm.  All five timing cases return the same conclusive
verdict in both arms.  All $406$ native and research source pins remain
unchanged, and there are no failed or unreturned calls.

The experiment uses CPU~5 affinity on a shared six-logical-CPU host,
with Python~3.13.14.  No competing heavy job from this task runs during
measurement.  Other repository jobs may remain active, so affinity is
not an exclusive CPU reservation.  The exact command, environment,
before/after process records and successful exit are retained in
\path{data/pachner-descent-environment.json}.

The following table gives median milliseconds for each arm, the median
of the three paired baseline/enabled time ratios, and the baseline
A/A ratio.  The paired ratio is computed from matched runs; it is not
the quotient of the two marginal medians.  Values above one favor the
enabled arm for that paired statistic.
\begin{center}\small
\begin{tabular}{lrrrr}
\toprule
Input & Baseline ms & Enabled ms & Paired ratio & A/A\\
\midrule
Circle & 11.075 & 10.973 & 1.0125 & 1.0238\\
Trefoil & 1964.002 & 2118.615 & 0.9200 & 0.9511\\
Figure eight & 6477.670 & 5607.515 & 1.1413 & 1.0063\\
$TT$ & 24118.429 & 23867.372 & 1.0270 & 0.9949\\
$T_{25}$ & 19226.410 & 11190.667 & 1.7181 & 1.0174\\
\bottomrule
\end{tabular}
\end{center}

The trefoil is an observed regression: all three paired ratios are
below one, despite fewer normal-search nodes and pivots.  The timings
charge the geometric search and proof overhead needed to obtain that
smaller query.  The figure-eight median favors the enabled policy, but
one of its three individual ratios is below one.  The $TT$ difference
is small on this host, and the $T_{25}$ control has the largest observed
paired improvement.  These differing outcomes are retained together;
the policy remains optional.

The circle executes no macro search and serves as a protocol control;
its small timing difference is not evidence of a geometric performance
gain.  Three paired samples and a single A/A pair per input do not
establish a population-level effect or an asymptotic scaling law.
The benchmark measures these complete fixed-policy recognition calls,
including replay and serialization.  Its ratios do not isolate the
cost of one Pachner move, prove that lower $h$ always gives a faster
run, or close the missing general quasi-polynomial bound.
